\documentclass[11pt]{article}
\usepackage[T1]{fontenc}
\usepackage{lmodern,amsmath,amssymb,amsthm,booktabs}
\usepackage[margin=1in]{geometry}
\usepackage[colorlinks=true,urlcolor=blue,citecolor=blue]{hyperref}
\newtheorem{theorem}{Theorem}
\newtheorem{proposition}[theorem]{Proposition}
\newcommand{\rem}{\operatorname{rem}}
\newcommand{\eqzero}{\operatorname{eq}_0}
\newcommand{\Gen}{\operatorname{gen}}
\newcommand{\Pick}{\operatorname{choose}}
\newcommand{\Set}{\operatorname{set}}
\newcommand{\AcceptCandidate}{\operatorname{accept}}
\newcommand{\Loop}{\operatorname{loop}}
\newcommand{\lean}[1]{\nolinkurl{#1}}
\makeatletter\g@addto@macro\UrlBreaks{\do\_}\makeatother
\emergencystretch=2em
\title{Fixed-arity nondeterministic remainder programs:\\expanded proof and executable checks}
\author{Research note accompanying the uniform-process construction}
\date{October 16, 2026}
\begin{document}
\maketitle

\begin{abstract}
The Euclidean-prefix construction gives a family of finite nondeterministic
recursive programs with arbitrarily small logarithmic remainder-call
coefficients. Each program uses only remainder and the zero test, and its
arity is fixed independently of the inputs. This does not give a single
sublogarithmic program. We give the program schema, correctness and cost
arguments, and an executable compilation into first-order syntax.
Lean checks the fixed-register recurrence, the asymptotic program bound,
and translations preserving remainder counts into finite first-order syntax
and a specialized two-stack recursive machine.
It also checks correctness of every terminating syntactic evaluation and
existence of a successful evaluation on every input.
\end{abstract}

\section{The result and its quantifiers}

We use the nondeterministic convention of Moschovakis: every convergent
computation must return the correct value, divergent computations are
disregarded, and remainder complexity is the minimum number of remainder
calls in a convergent computation \cite[\S2E, (2E-2), (2E-5)]{book}.
All logarithms below have base $2$.

\begin{theorem}\label{thm:family}
For each integer $p\geq1$, there is one finite nondeterministic recursive
program $P_p$ over $(\mathbb N,\rem,\eqzero)$ computing the gcd such that
\[
\limsup_{a\to\infty}\;
\max_{1\leq b\leq a}\frac{c_{P_p}(a,b)}{\log a}\leq\frac3p.
\]
Its generator has $W=4^p+5$ individual parameters; the largest auxiliary
arity is $W+1$, from the setters below. Inlining the setters, as in the
companion manuscript and in Lean, gives maximal arity $W$. In particular,
\[
\forall\varepsilon>0\ \exists P, B\ \forall a\geq b\geq1,
\qquad c_P(a,b)\leq\varepsilon\log a+B.
\]
\end{theorem}

The program may depend on $\varepsilon$, but not on $a,b$.
This is compatible with the Main Conjecture, whose lower-bound coefficient
may depend on the program \cite[p.~2]{book}.
No passage from this family to one sublogarithmic program is asserted.
The bound concerns remainder calls, not internal recursive-machine steps
or the time needed to find a successful branch.

\section{A finite program for each fixed width}

Fix $W\geq2$ and write $\vec x=(x_0,\ldots,x_{W-1})$.
Define the following finite family of equations:
\begin{align*}
\Pick(\vec x)&=x_i &&(0\leq i<W),\\
\Set_i(\vec x,z)&=\Gen(x_0,\ldots,x_{i-1},z,x_{i+1},\ldots,x_{W-1})
  &&(0\leq i<W),\\
\Gen(\vec x)&=\Pick(\vec x),\\
\Gen(\vec x)&=\Set_i(\vec x,
   \rem(\Pick(\vec x),\Pick(\vec x))) &&(0\leq i<W),\\
\Gen(\vec x)&=\Set_i(\vec x,\Gen(\vec x)) &&(0\leq i<W),\\
\Gen(\vec x)&=\Gen(\Pick(\vec x),\ldots,\Pick(\vec x)).
\end{align*}
There are exactly $W$ occurrences of $\Pick$ in the last equation.
The repeated definitions express finite nondeterministic choice.
Each index $i$ is a syntactic index: after fixing $W$, all the displayed
families are expanded into finitely many ordinary equations.
The symbols $W,i$ are not runtime data, and there is no array-access,
list-construction, or numeric-constant primitive.

The last equation copies or permutes any available parameters without a
remainder call. In the subcomputation equation
$\Set_i(\vec x,\Gen(\vec x))$, call-by-value evaluates the inner generator
once and retains the caller's parameters until its return. The result replaces
one parameter, and the outer generator continues.
The Lean presentation inlines the setters: for example, the remainder
alternative becomes
$\Gen(\vec x)=\Gen(x_0,\ldots,x_{i-1},
\rem(\Pick(\vec x),\Pick(\vec x)),x_{i+1},\ldots,x_{W-1})$.
The subcomputation alternative is inlined in the same way. Its fixed
signature has four internal symbols: choose, gen, accept, and loop.
The executable Python presentation retains the setters.
These are ordinary first-order calls; each recursive variable has fixed
arity \cite[\S\S2A, 2D]{book}. The recursion stack may grow, so $W$ is not a
bound on the total number of values in the entire computation.

Complete the program by setting $\Loop(z)=\Loop(z)$ and
\[
\AcceptCandidate(a,b,g)=
\begin{cases}
\Loop(g),&\eqzero(g),\\
g,&\eqzero(\rem(a,g))\ \text{and}\ \eqzero(\rem(b,g)),\\
\Loop(g),&\text{otherwise}.
\end{cases}
\]
The cases are nested conditionals: the zero test is performed before either
division, and the conjunction is expanded into two conditionals.
Finally set
\[
P_W(a,b)=
\begin{cases}
a,&\eqzero(b),\\
\AcceptCandidate(a,b,\Gen(a,b,a,\ldots,a)),&\text{otherwise}.
\end{cases}
\]
Only $\rem$ and $\eqzero$ are primitive symbols in this program.
On a successful computation with $b>0$, verification adds exactly two
remainder calls and the whole program makes four zero tests.

\section{Correctness and the operational interpretation}

\begin{proposition}
For every $W\geq2$, every convergent computation of $P_W(a,b)$ returns
$\gcd(a,b)$, and at least one such computation exists on every natural input.
Positive scaling preserves the generator's successful choice sequences and
remainder-call counts.
\end{proposition}
\begin{proof}
Induct on a finite generator evaluation. Every common divisor of its input
parameters divides its output: this is immediate for a selected parameter,
is preserved by remainder and copying, and passes through a recursive
subcomputation and its continuation.

Consequently $\gcd(a,b)$ divides every candidate $g$ generated at the root.
A successful check requires $g>0$, $g\mid a$, and $g\mid b$, so also
$g\mid\gcd(a,b)$. Thus $g=\gcd(a,b)$.
The separate zero-input branch is correct as well.

There is always a successful generator computation following Euclid.
With the current pair in the first two parameters, overwrite the first by
its remainder modulo the second, then swap and duplicate parameters to obtain
the next pair. When the second is zero, return the first.
The choice sequence is selected in the existence proof; no additional test
is introduced into the generator.
Finally $\rem(gx,gy)=g\rem(x,y)$ for $g>0$, and nonzero divisors remain
nonzero under scaling. Induction preserves the same computation and its cost.
\end{proof}

The Lean relation \lean{ShortGCD.RegisterProgram.Run} has four constructors:
return a parameter, overwrite one parameter by a remainder, route parameters,
or run a subcomputation and overwrite one parameter by its result.
These correspond to the four generator alternatives above. For a remainder
step, the two selector calls choose its operands; for a routing step, the $W$
selector calls choose its argument tuple. Internal selector calls make no
remainder calls.

The kernel-checked theorem \lean{Syntax.run_invokes} translates every register
run into an evaluation of the inlined first-order generator, preserving its
remainder count exactly. \lean{Syntax.accepts_eval} includes the verification
wrapper. Independently, \lean{Syntax.eval_head_correct} proves correctness of
\emph{every} terminating evaluation, and \lean{Syntax.eval_head_exists}
proves existence of a successful evaluation on every input. Thus correctness
is not restricted to the selected fast witnesses.

The syntax contains variables, remainder, zero-test conditionals, and calls
to a fixed finite signature; each symbol has a finite set of bodies.
\lean{Syntax.Eval} uses call-by-value evaluation trees, evaluating all arguments
before the chosen body and summing their remainder counts.
\lean{RegisterSmallStep} gives an explicit two-stack machine for this language:
it evaluates arguments right to left and has separate transitions for
zero tests and branching. Its closures encode substituted source terms,
not a new data type available to the program. The Python
compiler and interpreter are independently checked executable artifacts,
not verified implementations of the Lean semantics.

For these pure terms the evaluation trees correspond, with the same result
and remainder count, to the book's recursive machine. In the forward direction,
induct on a tree and serialize its argument evaluations right to left before
entering the chosen body. Transition locality embeds each evaluation in its
caller context \cite[Table~1 and Lemma~2D.1, pp.~75--76]{book}. A remainder
node adds one remainder call; a zero-test conditional adds none and evaluates
only its chosen branch. Conversely, decompose a terminating machine computation
into its completed argument computations and selected body or branch, and
induct on its length. Each selector occurrence makes its own choice, with no
shared mutable state. The setters in this note merely pass their evaluated
arguments to the generator, so inlining them also preserves results and
remainder counts. \lean{Syntax.Machine.eval_iff_steps} now checks the
tree-to-machine equivalence in both directions; \lean{Machine.Step.context}
checks transition locality. The machine specializes the book's transition
table to the fixed language and signature, rather than formalizing arbitrary
structures and vocabularies. The separate equivalence with the Python
presentation retaining setters remains a mathematical argument.

\section{The fixed-prefix cost bound}

We use the Euclidean-prefix and three-call contraction lemmas from the
companion manuscript \cite{companion}. Fix $p\geq1$, and put
\[
k=2^p,\qquad h=k^2+1=4^p+1,\qquad W=h+4.
\]
The entire prefix $r_0,\ldots,r_h$ fits in $h+1$ generator parameters.
Three further parameters suffice for the current modulus and two temporary
residues. Each contraction uses at most three remainder calls and can
overwrite these temporary parameters, leaving the prefix available.

After reaching a small modulus $m$, form $u=a\bmod m$ and $v=b\bmod m$
in the two temporary positions. Use a recursive generator subcomputation
to produce $d=\gcd(m,u)$, first routing its parameters to copies of $m,u$.
Its caller retains $v$. On return, replace the modulus parameter by $d$,
then route to copies of $d,v$ for the second gcd computation.
The identity
\[
\gcd(\gcd(m,a\bmod m),b\bmod m)=\gcd(m,\gcd(a,b))
\]
recovers the original gcd, since every generated modulus is divisible by it.
Only one result crosses each subcomputation boundary.
In the early-small-remainder case, route directly to the next Euclidean pair.
Thus the two recursive branches of the certificate proof are realizable in
one fixed generator of width $W$, with the same remainder counts.

For precision, let $T_W(n)$ be the maximum, over $a,b\leq2^n$, of the minimum
generator remainder count needed to produce $\gcd(a,b)$ from the initial
tuple $(a,b,a,\ldots,a)$. Euclid gives $T_W(n)=O(n)$, and copying gives symmetry.
Positive scaling lets us analyze coprime inputs and transport the resulting
computations to general inputs without extra calls.

Fix $0<\beta<1/2$ and set $\alpha=1-\beta/2$.
Use the early threshold $N^\alpha$ and the target modulus $N^\beta$, with
$N=2^n$. If the prefix is large, then
$2a/r_h<2N^{1-\alpha}<N^\beta$ for sufficiently large $n$.
The prefix therefore occupies at least $k+1$ residue classes modulo every
modulus used before stopping. A contraction decreases the modulus by a factor
greater than $k$, and $J=\lceil n/p\rceil$ contractions suffice.
It follows that
\[
T_W(n)\leq
\max\left\{
 T_W(\lceil\alpha n\rceil)+h+1,\quad
 2T_W(\lceil\beta n\rceil)+h+3\lceil n/p\rceil+2
\right\}.
\]
All thresholds and operand choices select witnesses among computations of
the same fixed program; $\beta$ is not a program parameter.

Write $L_W=\limsup_{n\to\infty}T_W(n)/n$, finite and nonnegative by Euclid.
Here $h$ is fixed, so dividing the recurrence by $n$ and taking upper limits
gives
\[
L_W\leq\max\{\alpha L_W,\ 2\beta L_W+3/p\}.
\]
Since $\alpha<1$, we obtain $(1-2\beta)L_W\leq3/p$.
Letting $\beta\downarrow0$ yields $L_W\leq3/p$.
Taking $n=\lceil\log a\rceil$ and adding the two verification calls proves
the first assertion of Theorem~\ref{thm:family}. For the second, choose
$p>3/\varepsilon$ and absorb the finitely many smaller inputs in $B$.

This argument uses a fixed $h$, unlike the sublogarithmic certificate
construction. The price for a smaller logarithmic coefficient is a larger
finite program. Neither the stated coefficient nor the parameter count is
claimed optimal.

\section{What still prevents the direct sublogarithmic compilation}

The original construction uses a prefix length growing with input size.
Replacing $W$ by a function of the input in the compiler would produce
different programs on different inputs. It would not prove the desired
single-program result.

Saving a growing prefix in nested calls also does not directly implement the
reservoir. A current function body can use its individual parameters; values
retained in older continuations are not an indexed store. Returning exposes
an older continuation but passes only the return value from the completed
call. For example, a recursive prefix builder that carries only the current
Euclidean pair leaves the innermost call without all earlier remainders as
parameters. The reservoir lemma cannot be applied there as though those
values were freely available.

Regenerating the whole prefix for every contraction is too expensive for
this construction: a positive prefix $r_0,\ldots,r_h$ requires $h-1$
remainder calls, so $J$ rebuilds cost $(h-1)J$ calls before the contractions.
For $h$ of order $\sqrt n$ and $J$ of order $n/\log n$, this is of order
$n^{3/2}/\log n$.

This diagnoses the direct implementations, not every possible recursive
program. A different organization of the computations might avoid the need
for a fully reusable growing prefix. No impossibility theorem for such an
organization has been proved here. In particular, the quantifier exchange
\[
\forall\varepsilon\ \exists P_\varepsilon
\qquad\hbox{to}\qquad
\exists P\ \forall\varepsilon
\]
remains unjustified.

\section{Executable checks and current status}

\lean{register_program.py} compiles a fixed-width instance to an AST containing
only individual variables, function calls, and conditionals. A static checker
checks arities and sorts and rejects extra primitives and numeric literals.
A separate evaluator implements call-by-value, with right-to-left argument
evaluation. It executes the compiled syntax and counts actual remainder calls.
The offline planner supplies nondeterministic choices; its unrestricted
arithmetic is not executed by the program.

The suite reuses each of the three compiled programs of widths $9,21,69$
across 138 cases, including zero, unordered, scaled, random, and Fibonacci
inputs. It additionally replays 129 positive-input cases after scaling both
inputs by 13 with exactly the same choices. The suite exercises 2,086
contractions and 101 recursive subcomputations, with nesting depth up to four.
Eight negative checks cover rejected candidates, zero divisors, malformed
choices, and invalid program syntax. A two-register instance separately
replays Euclid.

On one seeded 2,048-bit pair, whose ordinary Euclidean computation takes
1,188 remainder calls, the fixed programs give:
\begin{center}
\begin{tabular}{rrrr}
\toprule
$p$ & Generator parameters & Remainder calls & Internal calls\\
\midrule
1 & 9 & 1,036 & 4,476\\
2 & 21 & 713 & 3,564\\
3 & 69 & 634 & 4,920\\
\bottomrule
\end{tabular}
\end{center}
The remainder counts include final verification. These are successful
computations, not evidence of optimal costs or a deterministic method of
finding the choices. Fibonacci examples are also recorded, but should not
be taken as evidence of a new lower-bound separation: Pratt's program is
already known to be fast on that family \cite[Corollary~6C.7]{book}.

Run from the repository root:
\begin{verbatim}
lake env lean research/RegisterProgram.lean
python3 research/check_register_program.py
python3 -O research/check_register_program.py
python3 research/register_program.py --width 9
\end{verbatim}
The program-family bound is now kernel-checked as well as available for
independent mathematical scrutiny. The main declarations, all under
\lean{ShortGCD.RegisterProgram}, are:
\begin{itemize}
\item \lean{maxCost_step}: the fixed-register finite recurrence;
\item \lean{Fixed.limsup_profile_le}: the leading coefficient $3/p$;
\item \lean{Fixed.eventually_accepts_bound}: the bound for accepted runs,
including the final verification;
\item \lean{Syntax.eventually_eval_bound}: the same bound for first-order
program evaluations;
\item \lean{Syntax.arbitrarily_small_program}: the quantified
$\varepsilon\log a+B_\varepsilon$ statement with the program fixed before
the inputs are chosen.
\item \lean{Syntax.Machine.eval_iff_steps}: exact-cost equivalence between
evaluation trees and terminating small-step executions;
\item \lean{Syntax.Machine.arbitrarily_small_program}: the same quantified
bound directly for finite machine transition sequences.
\end{itemize}
The proof modules are included in the normal \lean{lake build} target.
The research entry point prints their axiom dependencies, recorded in
\lean{register-axioms.txt}; only the standard axioms \lean{propext},
\lean{Classical.choice}, and \lean{Quot.sound} occur.
The program-family theorem and its proof are included in Section~6 of the
companion manuscript; this note retains the expanded operational discussion
and executable checks. The accompanying \lean{model-novelty-audit.txt}
records the model comparison and the scope of the primary-literature search.
Obtaining one sublogarithmic program remains the larger open task.
The separate investigation \lean{single-program-investigation.pdf} tests a
fixed program that regenerates selected prefix operands, analyzes its
accounting limitation, and identifies a possible structural direction.

\begin{thebibliography}{2}
\bibitem{book}
Yiannis N. Moschovakis.
\emph{Abstract Recursion and Intrinsic Complexity}.
Lecture Notes in Logic 48, Cambridge University Press / ASL, 2019.
\href{https://doi.org/10.1017/9781108234238}{doi:10.1017/9781108234238}.
\bibitem{companion}
Jason Davies.
\emph{The Euclidean algorithm is not optimal among uniform processes}.
Companion repository manuscript, October 2026.
\end{thebibliography}
\end{document}
